\section{Catholic demonology: doctrine, inference, and restraint}

The Church's discipline makes sense only within a sharply asymmetrical
theology. God is uncreated goodness, the source of every creature and every
power creatures possess. Lateran IV teaches that the devil and other demons
were created good by nature but became evil by their own doing (const.~1,
\emph{Firmiter credimus}, DS 800). The Catechism identifies them as fallen
angels, describes their rejection of God as free and irrevocable, and insists
that Satan remains a creature whose finite action is permitted by providence
(CCC 391--395). Evil is therefore neither an eternal principle nor a second
sovereignty.

This is doctrine. Many details supplied by later demonologies are not. The
existence and fall of personal evil spirits, their opposition to God's reign,
their temptation of human beings, Christ's victory, and the final defeat of
evil belong to the Church's received confession (CCC 391--395, 550,
2850--2854). This study's authority audit found no defined status for elaborate
ranks or territorial assignments, non-scriptural names, causal explanations of
a particular misfortune, or pastoral taxonomies of ``infestation,''
``oppression,'' and ``obsession.'' Such terms may distinguish questions in a
particular school or pastoral setting; repetition does not make them dogma.

\subsection{Four source functions}

An account can repeat only true propositions and still mislead if it suppresses
their different functions and kinds of authority. The following four-part map
is not a rank order within the magisterium.
The first is the conciliar confession received in Lateran IV: the devil and
the other demons are creatures, good in their created nature and evil through
their own act. The second is authoritative catechetical synthesis. The
Catechism gathers the fall of the angels, their irrevocable choice, their
finite creaturely power, providence, Christ's victory, temptation, sacramentals,
human dignity, and Christian prayer into one account (CCC 391--395, 550,
1667--1673, 1700--1706, 2850--2854). The synthesis is more extensive than the
one conciliar sentence;
its propositions should not all be attributed to that sentence as though one
locus contained the whole later account.

The third level is school theology. Aquinas offers a disciplined analysis of
angelic knowledge, influence, temptation, and baptismal preparation. His
arguments clarify what a created intellect and will can and cannot do, but this
study does not turn every conclusion of that school into a defined proposition.
The fourth level is editorial synthesis. Terms such as natural,
preternatural, and supernatural, or a comparison among doctrine, authorization,
and diagnosis, can help prevent category mistakes. Their usefulness does not
give the editor authority to create a Catholic taxonomy.

These levels are related without being interchangeable. Doctrine supplies
boundaries within which theology reasons. Catechesis orders many doctrines for
teaching and Christian life. School theology asks further questions and
proposes arguments. Editorial synthesis explains why the distinctions matter
for this history. A claim does not become dogma because a theological school
is venerable, a pastoral term is widespread, or a contemporary author places
it beside a conciliar citation.

\begin{threestable}{Level}{What it contributes here}{What it cannot establish}
\textbf{Conciliar confession} & The created goodness and creaturely fall of
the devil and demons. & A complete demonology, pastoral taxonomy, diagnosis,
or ritual procedure.\\
\textbf{Catechetical synthesis} & The fall, finite creaturely power,
providence, Christ's victory, sacramentals, dignity, and prayer. & A finding
about a particular person or the permission required by law.\\
\textbf{School theology} & Reasoned distinctions concerning angelic knowledge,
influence, temptation, and a scholastic interpretation of inherited ministry.
& Defined doctrine merely by
being traditional, or a clinical account of symptoms.\\
\textbf{Editorial synthesis} & A transparent map of categories and evidentiary
limits. & A new doctrine, universal Catholic terminology, or authority over a
case.\\
\end{threestable}

\subsection{Creation, fall, and the asymmetry of evil}

The created-good and self-made-evil clauses do more work than deny a rival god.
They separate nature from moral defection. Demons are not evil because
creaturehood, spirit, intelligence, or power is evil. Their evil is not a
substance added by a second creator. The doctrinal grammar therefore blocks
both dualism and a fearful sacralization of evil. Whatever power a creature
has remains received, finite, and subject to providence (CCC 391--395).

This asymmetry also governs historical interpretation. A source may devote
many lines to hostile spirits because its literary subject requires it; that
verbal prominence cannot make evil metaphysically coordinate with God.
Conversely, a sparse treatment does not prove that an author or community
denied personal evil. Genre, purpose, and authority must be asked before the
length or drama of a passage is treated as theological weight.

The limit of the conciliar control is equally important. The exact witness
used here is the Holy See's later reception and quotation of
\emph{Firmiter credimus}, with its DS and conciliar locus, not a new collation
of the 1215 act. It controls the attributed proposition. It does not by itself
control the complete medieval context, the constitution's manuscript history,
or every later interpretation. The Catechism supplies a separately identified
authoritative synthesis; it is not silently substituted for the conciliar
artifact.

\subsection{The Christological center}

The Gospel controversies about expulsion are arguments about Christ and the
kingdom before they are material for a later theory of cases. In the Synoptic
stronger-man sayings, the decisive question is the source and meaning of
Jesus' authority (Mk 3:22--27; Mt 12:22--32; Lk 11:14--22). The Catechism's
summary makes the same center explicit: Christ's expulsions manifest the
arrival of God's reign and his victory over the ruler of this world (CCC 550).
The theological subject is not an autonomous map of demonic capacities but
the identity and mission of Christ.

That center disciplines later reception in two directions. It prevents a
naturalistic historical method from deciding in advance that the narratives
can have no theological referent; the texts must first be heard as the
canonical witnesses they are. It also prevents a Christian reader from
treating them as a portable diagnostic or procedural manual. Narrative action,
controversy, proclamation, and later ecclesial discipline are distinct genres.
The Church's confession of Christ's victory does not supply a checklist by
which a reader can assign an unseen cause to another person's suffering.

The petitions concerning deliverance from evil in Christian prayer likewise
place spiritual conflict within dependence on the Father, forgiveness,
temptation, perseverance, and final hope (CCC 2846--2854). This broad horizon
corrects a reception that makes extraordinary affliction the normal measure of
Christian seriousness. Ordinary temptation and the struggle for faithful
consent concern every Christian; a reserved rite does not.

\subsection{Dogma does not function as a case archive}

The doctrinal propositions and the history of alleged cases answer different
questions. Lateran IV identifies the created status and moral fall of demons.
The Catechism orders that confession within creation, providence, Christology,
moral life, sacramentals, and prayer. Neither source supplies a census of
affliction, a dossier of verified events, or an empirical rate. A dogmatic
proposition can be true without making every narrative offered under its
vocabulary historically reliable.

The converse also follows. Failure to verify a particular report does not
disprove the doctrinal proposition to which its narrator appeals. Historical
criticism may find dependence, literary shaping, missing records, or
contradiction in a case tradition. Those findings govern that tradition. They
do not become a laboratory judgment about angels. This separation protects
both inquiries: theology is not made hostage to sensational accounts, and
history is not required to certify an event because the metaphysical category
is Catholic.

Reception often obscures the separation by moving authority sideways. A
catechetical statement is cited as though it authenticated a memoir; the
memoir is then cited as though it supplied empirical confirmation of the
catechesis. The two references appear mutually supporting while neither
performs the task assigned to it. A disciplined account instead asks of each
claim whether its support is doctrinal, juridical, historical, pastoral, or
clinical, and then stops at that source class's limit.

\subsection{Natural, preternatural, and supernatural}

For the bounded purposes of this study, these school-theological terms mark
relations, not three kinds of being. The definitions that follow are an
editorial synthesis used to keep kinds of claim distinct; the three-word
scheme is not presented as a definition of doctrine.
\emph{Natural} refers to a creature and to powers and effects proportioned to a
created nature. Catholic writers may call an alleged effect
\emph{preternatural} when they attribute it to a created spiritual power beyond
the ordinary capacities of embodied human agents. An angel or demon is
nevertheless a creature with an angelic nature, not a ``preternatural
substance.'' \emph{Supernatural} refers to God and to the order exceeding every
created nature: uncreated divine life and created gifts of grace by which
creatures participate in it.

These are metaphysical classifications, not diagnostic test results. A
surprising event does not become preternatural because its natural explanation
is not yet known. ``Unexplained'' describes an epistemic state; it does not
name a cause. Conversely, Catholic theology does not require that the
preternatural be redefined as an unknown natural mechanism. It permits the
question while refusing a shortcut from possibility to case judgment.

The same distinction governs medicine. A clinician can diagnose, treat, or
exclude conditions within professional competence. Clinical evidence cannot
demonstrate that no created spirit exists; neither can a theological belief
diagnose epilepsy, psychosis, dissociation, delirium, intoxication, or trauma.
At the direct intersection---a living person's symptoms and safety---current
U.S. pastoral guidance calls for competent medical and psychological or
psychiatric assessment. This is a safeguarding discipline, not a clinical test
for a spiritual cause (CCC 1673; USCCB, ``Exorcism'').

\subsection{Temptation, sin, and responsibility}

The ordinary Christian doctrine of temptation is broader and more important
than extraordinary claims. Being tempted is distinct from consenting:
personal sin results from morally attributable consent. Responsibility follows
voluntariness, so ignorance, duress, fear, habit, and psychological or social
factors may diminish or nullify imputability; mortal sin requires grave matter,
full knowledge, and deliberate consent (CCC 1734--1735, 1857, 1859,
2846--2847). Neither grave sin nor persistent habit proves direct demonic
action. Persons remain moral subjects even when freedom or knowledge is
impaired. Attributing harmful conduct to possession can displace
accountability, blame a harmed person, stigmatize illness, or conceal abuse.

In Aquinas's scholastic analysis, demons neither know secret thoughts as they
exist in intellect and will nor move the will inwardly or by necessity. They
may infer dispositions from outward effects, persuade through an apparent
good, rouse passions, and affect imagination or sense through bounded bodily
causality; the will remains able to consent or resist (ST I, q.~57, a.~4;
q.~111, aa.~2--4; q.~114, aa.~2--3). Aquinas also denies that every sin comes
from direct demonic instigation. This metaphysical account is neither defined
clinical mechanism nor a warrant for attributing a particular symptom.

\subsection{A school account and its limits}

Those distinctions restrain opposite errors. One error makes a demon the
author of everything a distressed or sinful person does, dissolving human
agency and the responsibility of other human actors. The other treats every
theological claim about creaturely spiritual agency as though it were a
primitive rival to medicine. Aquinas is answering metaphysical and moral
questions, not supplying a differential diagnosis. His account cannot identify
the cause of speech, movement, memory, fear, intrusive thought, or altered
behavior in a living person.

Nor should the limits be hidden beneath the authority of his name. This study
does not claim that his psychology is a clinical mechanism, that every Catholic
school uses identical terminology, or that his thirteenth-century discussion
describes a present diocesan process. In the history of offices, his treatment
of the exorcist within baptismal preparation witnesses one scholastic
interpretation of inherited ministry (ST III, q.~71, aa.~2--4). It cannot
collapse baptismal exorcism, a former minor order, and the current reserved act
into one timeless office.

\subsection{Possession and the integrity of the person}

In Catholic pastoral usage, alleged possession concerns extraordinary
affliction attributed to demonic power; it never means that a human being has
become a demon. The person remains the subject of human dignity and care (CCC
1700--1706). An allegation does not authorize treating every utterance during
distress as a spirit's speech, ignoring disclosures of pain or abuse, or
relabeling refusal as demonic resistance. Capacity must be assessed, consent
sought wherever possible, and applicable safeguarding and legal protections
remain operative (CCC 1673; USCCB, ``Exorcism'').

The rite is directed toward liberation through Christ and the prayer of the
Church. It is not punishment, interrogation, evidence extraction, or an ordeal
designed to force a reaction. The person seeking or referred for care is never
the enemy. This Christological boundary is more fundamental than any proposed
sign: Christ has no equal rival, and the Church's ministry seeks liberation
rather than spectacle or domination (CCC 550, 1673).

\subsection{Doctrine, sacramental, and authorization}

Three affirmations that often appear together answer different questions.
First, doctrine identifies the creaturely reality of evil spirits, Christ's
victory, and the dignity of the afflicted person. Second, the Church classifies
exorcism among the sacramentals: sacred signs obtain principally spiritual
effects through the Church's intercession (CIC c.~1166; CCC 1667, 1673).
Third, Latin law reserves licit exorcism over the possessed to a qualified
priest who has obtained the local ordinary's particular and express permission
(CIC c.~1172). Doctrine supplies neither a diagnosis nor an appointment; the
definition of a sacramental supplies neither permission nor jurisdiction.

The word \emph{exorcist} consequently requires a dated context. It may denote a
holder of an ancient or medieval ecclesial grade, the minor order discussed in
scholastic and later Latin sources, or a priest permitted to undertake the
reserved ministry under current Latin law. Those are not interchangeable
offices. Aquinas's placement of exorcists among ministers assisting the priest
in baptismal preparation (ST III, q.~71, a.~4) does not confer the permission
required by canon 1172, describe a current diocesan post, or prove habitual
practice of major exorcism.

Nor does permission settle the factual judgment that precedes ritual action.
The Catechism's statement that major exorcism may be performed only by a priest
with episcopal permission gives catechetical orientation (CCC 1673); canon 1172
states the legal requirement as particular and express permission from the
local ordinary. Neither proposition turns authorization into proof that a
person is possessed. The office or administrative arrangement by which a
diocese receives referrals, the authority competent to give permission, and
the discernment of a particular presentation remain distinct.

The resulting category controls are concise: belief does not establish a
finding in a case; permission does not establish the alleged cause; clinical
diagnosis does not adjudicate every metaphysical proposition; and possession
of a book does not confer office.

\subsection{Angels, saints, sacramentals, and magic}

Sacramentals are sacred signs instituted by the Church; through her
intercession they signify and obtain effects, especially spiritual ones. They
do not confer grace in the manner of the sacraments, but prepare persons to
receive grace and dispose them to cooperate with it (CIC c.~1166; CCC 1667,
1670). Their use is not a mechanism under the user's control: attributing
efficacy to mere external performance apart from the dispositions it calls for
is superstition (CCC 2111). Blessed objects are not talismans and sacred words
are not incantations. The materiality of a sign does not make it magic. The
decisive contrasts are petition rather than control, ecclesial worship rather
than self-authorized technique, and trust in God rather than an attempt to tame
occult powers or place them at one's service (CCC 2116--2117).

Invoking angels or saints within received Christian prayer differs from
seeking secret names, bargaining with spirits, divination, necromancy, or
coercive technique (CCC 2116--2117). Curiosity about a demon's identity has no
pastoral claim over the person's dignity. The CDF's letter of 29 September
1985 forbids direct questioning aimed at learning a demon's identity in the
unauthorized prayer gatherings it addresses. That precise restriction should
not be expanded into a description of an authorized rite. Historical ritual
questions are not reader-facing instructions.

\begin{studybox}{Editorial synthesis: the doctrinal center}
Christian demonology is a small province within creation, providence,
Christology, and redemption. Its central propositions are that evil is
creaturely and parasitic, Christ is victorious, human dignity remains, and the
Church's prayer depends on God. A system that produces fascination, fatalism,
fear of ordinary medicine, or contempt for an afflicted person has lost that
center even when it uses traditional vocabulary.
\end{studybox}
